%---------------------------------------------------------------------------%
%->> Section 2.1: 微服务故障诊断基础概念
%---------------------------------------------------------------------------%

\section{微服务故障诊断概述}

微服务架构将复杂业务拆分为多个可独立部署的服务单元，各服务围绕特定业务能力运行，并通过 HTTP 或 RPC 协同处理请求。这种组织方式提升了系统扩展性和迭代效率，也改变了故障的形成方式。故障不再局限于单一模块内部，还可能沿服务调用、资源竞争或基础设施依赖向外扩散，表现为跨服务、跨层级的异常。这种分布式特性使得故障的可观测性和可追溯性都面临新的挑战。

微服务故障诊断的难点主要来自三个方面。第一，业务逻辑被拆分到多个服务及其支撑组件中，根因容易隐藏在下游依赖内部，导致症状与根因在空间上出现分离。第二，系统通常运行在容器编排环境中，实例会随调度和负载持续变化，静态拓扑难以完整反映运行时依赖。第三，一次请求往往穿过网关、业务服务、缓存、消息队列和数据库等多个环节，任一环节异常都可能沿调用链向上传播，并在入口侧形成更明显的外部症状，从而掩盖真实根因的位置。

以电商系统为例，用户请求经由网关进入系统后，通常要经过订单、支付、库存等多个服务以及数据库、缓存等支撑组件。若支付数据库连接池耗尽，异常会先影响支付服务，再沿调用关系传递至订单服务和网关，最终表现为页面超时或请求失败。可见症状与真实根因之间往往隔着若干服务和依赖环节，这正是微服务场景与单体系统诊断的重要区别。这种传播链条的存在要求诊断方法必须具备回溯能力。

根因分析（Root Cause Analysis，RCA）的目标，是从外部异常现象出发，结合运行证据与逻辑推断识别导致故障的真实原因。与针对表面症状的临时处置不同，根因分析关注故障为何产生、异常如何传播，以及究竟哪个组件或关系构成根因。准确的根因定位能够为后续修复提供依据，并降低同类故障重复发生的概率，从而提升系统整体可靠性。

故障定位任务可形式化表述如下。设整个微服务系统由一组服务节点$S=\{S_1, S_2, \ldots, S_m\}$及其调用关系$E \subseteq S \times S$构成。每个服务节点$S_i$会持续产生多模态监控数据$D_i = \{T_i, M_i, L_i\}$，分别对应调用链路、指标和日志。在某一时刻$t$，系统触发异常检测模块，得到候选异常告警集合$A_t = \{a_1, a_2, \ldots, a_k\}$。根因定位函数可定义为$\theta: A_t \rightarrow S \cup E$。其中，$S_i$表示被判定为根因的故障服务节点，$e_{ij} = (S_i, S_j)$表示被判定为根因的故障调用链路。也就是说，给定异常告警集合$A_t$后，模型需要输出最可能导致系统整体异常的故障位置。

微服务环境下的根因分析还面临几类共性挑战。故障传播路径往往较长，相同症状背后可能对应不同根因，导致诊断存在多解性；系统运行状态持续变化，静态拓扑难以完整描述运行时依赖；可用于监督学习的故障实例相对稀缺，监督信号有限；指标、日志和调用链路之间又存在明显异构性。如何在保留各模态特征的同时建立稳定关联，是微服务故障诊断的核心问题。
